% 自编教学补充；阅读来源见本章 README。
\begin{frame}[fragile]{2.3A 直方图：每个柱子数什么}
\small 合成成绩用于解释分箱：每根柱子数一个分数区间内的学生人数。它与“每班一根柱子”的类别柱状图不同。
\medskip
\begin{lstlisting}
exam_scores = [55, 62, 68, 71, 75, 78, 82, 85, 91, 96]
fig, ax = plt.subplots(figsize=(6, 3.5))
ax.hist(exam_scores, bins=[50, 60, 70, 80, 90, 100],
        edgecolor="white", color="#2B6A6C")
ax.set(title="合成成绩的分布", xlabel="成绩（分）", ylabel="人数")
fig.tight_layout()
plt.show()
\end{lstlisting}
\end{frame}

\begin{frame}[fragile]{2.3B 散点图：一行观测对应一个点}
\small 每个点代表一名虚构学生。横轴学习时间、纵轴成绩；即使有正向关系，也不能仅凭这张图认定学习时间造成成绩差异。
\medskip
\begin{lstlisting}
study_hours = [1, 2, 2, 3, 4, 4, 5, 6, 6, 7]
fig, ax = plt.subplots(figsize=(6, 3.5))
ax.scatter(study_hours, exam_scores, alpha=0.8, color="#B40D49")
ax.set(title="学习时间与成绩（合成）", xlabel="学习时间（小时）", ylabel="成绩（分）")
ax.grid(alpha=0.2)
fig.tight_layout()
plt.show()
\end{lstlisting}
\end{frame}
